% GENERATED FILE. Edit canonical lessons, then run npm run generate:books.
% canonical-source-hash: fnv1a64:ea3c15022307b04d
% canonical-lessons: FR-C09-printemps, FR-C09-ete, FR-C09-automne, FR-C09-hiver, FR-C09-en-saison, FR-C09-saisons-practice

\chapter{The Four Seasons}
\label{ch:seasons}
\hlchaptermodality{11}

\begin{chapteropening}
\textbf{By the end of this chapter:} \emph{I can name the four seasons in French and say that something happens in one of them.}

Four small Latin poems about heat, sleep and first things. printemps is prīmum tempus, \textquotedblleft{}the first time\textquotedblright{}; été is aestās, \textquotedblleft{}heat\textquotedblright{}; hiver is hībernum, the wintry season that gave English hibernate. And three seasons take en while spring takes au, for a reason that is not about seasons at all.
\end{chapteropening}

\section[le printemps]{le printemps — spring}

\label{lesson:FR-C09-printemps}



\begin{quote}
Twelve months name gods, emperors and numbers. The four seasons are something else entirely: four small Latin poems about \textbf{heat, sleep and first things}. Here is the first thing.
\end{quote}

\begin{sounds}
\begin{itemize}
\raggedright
  \item \textbf{le printemps} = \emph{luh pran-TAHN}. \textbf{Two nasal vowels in one word}: the \textbf{in} of \emph{prin-} and the \textbf{em} of \emph{-temps}.
  \item Everything after that second nasal is silent. The \textbf{-ps} on the end is not said at all.
\end{itemize}
\end{sounds}

\begin{cousinweb}
\textbf{printemps} is two words fused: \textbf{prīmum tempus}, \textquotedblleft{}the \textbf{first time}.\textquotedblright{}

\begin{quote}
\textbf{prin-} = \emph{prīmum}, \textquotedblleft{}first\textquotedblright{} — English \textbf{prime}, \textbf{primary}, \textbf{premier} \textbf{-temps} = \emph{tempus}, \textquotedblleft{}time, season\textquotedblright{} — English \textbf{temporal}, \textbf{tempo}
\end{quote}

So spring is the year's \textbf{prime time}, and the French say so every time they name it. Older French used \textbf{primevère}, from the same \emph{prima vera}, \textquotedblleft{}first spring\textquotedblright{} — which Italian keeps to this day as \textbf{primavera}.
\end{cousinweb}

\subsection*{Your turn}
Say these aloud:

\begin{itemize}
\raggedright
  \item \textquotedblleft{}le printemps\textquotedblright{} — \emph{luh pran-TAHN}, both vowels through the nose
  \item it split — \textquotedblleft{}prīmum tempus: the first time\textquotedblright{}
  \item \textquotedblleft{}décembre, printemps\textquotedblright{} — the year ending, then the year beginning
  \item \textquotedblleft{}printemps\textquotedblright{} then Italian \textquotedblleft{}primavera\textquotedblright{} — two names for one idea
\end{itemize}

\begin{tcolorbox}[breakable,colback=teal!4,colframe=teal!35!black,title={Before you move on}]
How do you say spring? (\textbf{le printemps}, \emph{luh pran-TAHN}.) What are its two halves? (\emph{prīmum}, \textquotedblleft{}\textbf{first},\textquotedblright{} and \emph{tempus}, \textquotedblleft{}\textbf{time}.\textquotedblright{}) How many nasal vowels does it hold? (\textbf{Two}.) What does Italian call the same season? (\textbf{Primavera} — the same \textquotedblleft{}first spring.\textquotedblright{}) Next: the hot one.
\end{tcolorbox}
\section[l'été]{l'été — summer}

\label{lesson:FR-C09-ete}



\begin{quote}
Spring was named for its \textbf{place} in the year. Summer is named for what it \textbf{feels like}.
\end{quote}

\begin{sounds}
\begin{itemize}
\raggedright
  \item \textbf{été} = \emph{ay-TAY}. Both \textbf{é} are the bright, narrow \emph{ay} you met in \textbf{février}.
  \item \texttt{elision} — the article does not stay \emph{le}. Before a vowel sound it drops its own vowel and clamps on: \emph{le été} is impossible, so it becomes \textbf{l'été}, said as one word, \emph{lay-TAY}.
\end{itemize}
\end{sounds}

\begin{cousinweb}
\textbf{été} is Latin \textbf{aestās}, and \emph{aestās} means \textquotedblleft{}\textbf{heat}.\textquotedblright{}

Not \textquotedblleft{}the third season\textquotedblright{} or \textquotedblleft{}the warm months\textquotedblright{} — just \emph{heat}. Summer, in Latin and so in French, is named for the thing itself. English keeps the same root in \textbf{estival}, \textquotedblleft{}of summer,\textquotedblright{} and in \textbf{aestivate}, which is what an animal does when it sleeps through the hot season rather than the cold one.
\end{cousinweb}

\subsection*{Your turn}
Say these aloud:

\begin{itemize}
\raggedright
  \item \textquotedblleft{}l'été\textquotedblright{} — \emph{lay-TAY}, one word, the article clamped on
  \item \textquotedblleft{}le printemps, l'été\textquotedblright{} — hear the article change shape
  \item the meaning — \textquotedblleft{}aestās: heat\textquotedblright{}
\end{itemize}

\begin{tcolorbox}[breakable,colback=teal!4,colframe=teal!35!black,title={Before you move on}]
How do you say summer? (\textbf{l'été}, \emph{lay-TAY}.) Why is it \emph{l'} and not \emph{le}? (The word starts with a \textbf{vowel}, so the article elides.) What does the Latin \emph{aestās} mean? (\textbf{Heat}.) Next: the one English borrowed whole.
\end{tcolorbox}
\section[l'automne]{l'automne — autumn}

\label{lesson:FR-C09-automne}



\begin{quote}
Three of the four seasons are a Latin idea reshaped by French. This one is a Latin word that French barely touched, and then handed to English intact.
\end{quote}

\begin{sounds}
\begin{itemize}
\raggedright
  \item \texttt{diphthong-au} — the pair \textbf{au} is a single \emph{o} sound, closer to English \emph{taught} than to \emph{cow}. It never sounds like \emph{ow}.
  \item \textbf{automne} = \emph{o-TON}. Here is the oddity worth knowing: the \textbf{m} is written and the \textbf{n} is not sounded — the spelling keeps a letter the mouth dropped.
  \item With elision: \textbf{l'automne} = \emph{lo-TON}.
\end{itemize}
\end{sounds}

\begin{cousinweb}
\textbf{automne} is Latin \textbf{autumnus}, and unlike the other three seasons nobody is sure what \emph{autumnus} originally meant. It may not even be Latin in origin.

What is certain is where it went. English took \textbf{autumn} from French, spelling and silent letter and all — which is why English writes an \emph{n} it does not say in exactly the same place. Two languages carrying the identical silent letter, because one copied the other's spelling.

American English mostly prefers \textbf{fall}, a Germanic word for the leaves coming down; \textbf{autumn} is the French one.
\end{cousinweb}

\subsection*{Your turn}
Say these aloud:

\begin{itemize}
\raggedright
  \item \textquotedblleft{}l'automne\textquotedblright{} — \emph{lo-TON}, the \emph{au} a single \emph{o}
  \item \textquotedblleft{}l'été, l'automne\textquotedblright{} — both elide
  \item French \textquotedblleft{}automne\textquotedblright{} then English \textquotedblleft{}autumn\textquotedblright{} — the same silent letter twice
\end{itemize}

\begin{tcolorbox}[breakable,colback=teal!4,colframe=teal!35!black,title={Before you move on}]
How do you say autumn? (\textbf{l'automne}, \emph{lo-TON}.) What sound does \textbf{au} make? (A single \textbf{o}.) Which written letter is not pronounced, and does English do the same? (The \textbf{n} — and \textbf{yes}, English \emph{autumn} copied the spelling from French.) Next: the cold one, and a bear.
\end{tcolorbox}
\section[l'hiver]{l'hiver — winter}

\label{lesson:FR-C09-hiver}



\begin{quote}
The last season, and it shares a root with the deepest sleep in the animal world.
\end{quote}

\begin{sounds}
\begin{itemize}
\raggedright
  \item \texttt{silent-h} — the \textbf{h} is silent, exactly as in \textbf{heure}. And because a silent \emph{h} is no obstacle, the article elides straight through it: \textbf{l'hiver}, \emph{lee-VAIR}.
  \item The \textbf{r} at the end is the throat-\textbf{r}, and it \textbf{is} sounded.
\end{itemize}
\end{sounds}

\begin{cousinweb}
\textbf{hiver} is Latin \textbf{hībernum}, short for \emph{hībernum tempus} — \textquotedblleft{}the \textbf{wintry} season.\textquotedblright{} Like \emph{printemps}, it is a \emph{tempus} phrase with the \emph{tempus} worn off.

Now follow that root sideways. English \textbf{hibernate} is \emph{hībernāre}, and it means, with no metaphor at all, \textquotedblleft{}\textbf{to spend the winter}.\textquotedblright{} A bear hibernating is a bear \emph{wintering}.

So the same Latin word gave French a season and English a kind of sleep — one naming the time, the other naming what you do in it.
\end{cousinweb}

\subsection*{Your turn}
Say these aloud:

\begin{itemize}
\raggedright
  \item \textquotedblleft{}l'hiver\textquotedblright{} — \emph{lee-VAIR}, the \emph{h} silent and the article through it
  \item \textquotedblleft{}heure, hiver\textquotedblright{} — two silent \emph{h}, two elisions
  \item the pair — \textquotedblleft{}hiver and hibernate: the season and the sleeping\textquotedblright{}
\end{itemize}

\begin{tcolorbox}[breakable,colback=teal!4,colframe=teal!35!black,title={Before you move on}]
How do you say winter? (\textbf{l'hiver}, \emph{lee-VAIR}.) Why does the article elide in front of an \textbf{h}? (The \textbf{h} is silent, so the word begins with a vowel sound.) What does English \textbf{hibernate} literally mean? (\textquotedblleft{}\textbf{To spend the winter}.\textquotedblright{}) Next: how to say you are \emph{in} a season.
\end{tcolorbox}
\section[au printemps, en été]{au printemps, en été — saying you are \emph{in} a season}

\label{lesson:FR-C09-en-saison}



\begin{quote}
You have the four seasons. Now put something \emph{in} one — and meet a rule with exactly one exception, which turns out not to be an exception at all.
\end{quote}

\begin{grammarlens}[title={Grammar Lens: three take en, one takes au}]
Three of the four use \textbf{en}:

\begin{quote}
\textbf{en été} · \textbf{en automne} · \textbf{en hiver}
\end{quote}

And spring does not:

\begin{quote}
\textbf{au printemps}
\end{quote}

Learn it as a list and it is an oddity to memorise. Look at the words and it stops being odd. \textbf{Été}, \textbf{automne} and \textbf{hiver} all begin with a \textbf{vowel sound} — that is why each of them elides its article to \emph{l'}. \textbf{Printemps} begins with a consonant, and it is the only one that does.

So the split is not about seasons at all. It is the \textbf{same vowel-or-consonant question} that decided \emph{le} against \emph{l'}, asked one more time about a different little word. French keeps asking it, and once you hear it you stop needing the list.
\end{grammarlens}

\subsection*{Your turn}
Say these aloud:

\begin{itemize}
\raggedright
  \item the three — \textquotedblleft{}en été, en automne, en hiver\textquotedblright{}
  \item the one — \textquotedblleft{}au printemps\textquotedblright{}
  \item the reason aloud — \textquotedblleft{}three start with a vowel, one starts with a consonant\textquotedblright{}
\end{itemize}

\emph{Twice through:} \textquotedblleft{}au printemps, en été, en automne, en hiver\textquotedblright{}

\begin{tcolorbox}[breakable,colback=teal!4,colframe=teal!35!black,title={Before you move on}]
How do you say \textquotedblleft{}in summer\textquotedblright{}? (\textbf{en été}.) In winter? (\textbf{en hiver}.) In spring? (\textbf{au printemps} — not \emph{en}.) And \textbf{why} is spring different? (It is the only season starting with a \textbf{consonant}; the other three begin with a vowel sound, which is also why they elide to \emph{l'}.)
\end{tcolorbox}
\section[Practice]{Practice — the turning year}

\label{lesson:FR-C09-saisons-practice}



\begin{quote}
Four seasons, four small Latin poems. Run them, then put a month inside each one.
\end{quote}

\subsection*{Your turn: the four, in order}
Read down the left, then cover the right and read back:

\noindent
\begin{tabularx}{\linewidth}{@{}>{\raggedright\arraybackslash}X>{\raggedright\arraybackslash}X@{}}
\toprule
\textbf{French} & \textbf{English} \\
\midrule
\textbf{le printemps} & spring \\
\textbf{l'été} & summer \\
\textbf{l'automne} & autumn \\
\textbf{l'hiver} & winter \\
\bottomrule
\end{tabularx}

\emph{Twice through:} the four seasons, in order

Notice which article each one carries. \textbf{Three elide to \emph{l'}} and one does not, and that single fact is about to do a second job.

\subsection*{Your turn: in a season}
\begin{itemize}
\raggedright
  \item \textbf{au printemps} — the consonant one
  \item \textbf{en été} · \textbf{en automne} · \textbf{en hiver} — the three vowel ones
\end{itemize}

Now join what you learned last chapter to what you learned this one. Say each month inside its season:

\begin{quote}
\textbf{en hiver} — \emph{janvier, février} · \textbf{au printemps} — \emph{mars, avril, mai} \textbf{en été} — \emph{juin, juillet, août} · \textbf{en automne} — \emph{septembre, octobre}
\end{quote}

\begin{grammarlens}[title={What you've built this chapter}]
\begin{itemize}
\raggedright
  \item \textbf{four season names}, each one a Latin sentence worn down to a word: \emph{printemps} is \textquotedblleft{}the \textbf{first time},\textquotedblright{} \emph{été} is \textquotedblleft{}\textbf{heat},\textquotedblright{} \emph{hiver} is \textquotedblleft{}the \textbf{wintry} season,\textquotedblright{} and \emph{automne} is a word so old nobody is sure of it.
  \item \textbf{two of them still pay out in English} — \emph{estival} from \emph{aestās}, and \textbf{hibernate}, which means, with no metaphor, \textquotedblleft{}to spend the winter.\textquotedblright{}
  \item \textbf{au against en}, and the reason behind it: \textbf{printemps} is the only season beginning with a consonant. It is the same vowel-or-consonant question that decided \emph{le} against \emph{l'}, asked about a different word.
\end{itemize}
\end{grammarlens}

\begin{tcolorbox}[breakable,colback=teal!4,colframe=teal!35!black,title={Before you move on}]
Give the four seasons. Which three take \textbf{en}, and which takes \textbf{au}? (\emph{En été, en automne, en hiver}; \textbf{au printemps}.) Why? (Spring is the only one starting with a \textbf{consonant}.) What does \emph{printemps} literally mean? (\textquotedblleft{}The \textbf{first time}.\textquotedblright{}) Which English word means \textquotedblleft{}to spend the winter\textquotedblright{}? (\textbf{Hibernate} — from the same \emph{hībernum} as \emph{hiver}.)
\end{tcolorbox}
